\section{Conclusion}
We presented a SAT pipeline for minimum-span $L(h,k)$-labeling of simple
undirected graphs, with order-encoding equivalence, a span-search bound
invariant and conditions for valid symmetry breaking. The primary v2
comparison evaluates 117 fixed graph/parameter instances with three
repetitions and a 30-second limit per backend. SAT and assignment ILP
establish optimality in all three runs on \ProgressSATStable{} and
\ProgressILPStable{} instances, respectively. On the 21 initially open
instances, SAT has better upper bounds in \ProgressOpenSATUB{} of
\ProgressOpenPairs{} matched pairs, while ILP has stronger lower bounds
in \ProgressOpenILPLB{} pairs. These findings describe complementary
solution and proof behavior in this cohort, not general equivalence.

Table~\ref{tab:gap-coverage} consolidates the available evidence:
\ConstructionClosed{} of \ConstructionInstances{} intervals are closed
and \ConstructionOpen{} remain open. Square-clique bounds certify two
additional optima from saved witnesses, and a sufficient color-class
ordering condition supplies matching witnesses for three more instances.
In total, \ConstructionCliqueCertified{} optima can be checked directly
from a clique and a labeling, without relying on reported solver lower
bounds. This subset overlaps the closed total. The new certificates are
instance-specific and do not alter any original timing or status.

The separate 50-tree $L(3,2)$ study and archived $L(2,1)$ symmetry
experiments provide additional validation with their own protocols.
Neither finite agreement with reference formulas nor solver agreement
proves a new formula over an entire graph family. Existing SAT and ILP
labeling methods also preclude a first-use claim for these approaches.
The contribution is a documented exact implementation, correctness
arguments and reproducible evidence about its scope and limits.

A controlled direct-versus-order ablation and independently checkable
UNSAT traces could extend this work. GA remains a proposed separate
upper-bound baseline; no present result establishes superiority or
inferiority relative to it. Resolving the remaining intervals is a
possible extension rather than a prerequisite for reporting the
established results.
